а) Пусть $\alpha, \beta$ -- сопряжённые над полем $F$ элементы, 
$\beta \in F(\alpha)$. Тогда существует единственный автоморфизм 
$F(\alpha) \to F(\alpha)$, сохраняющий $F$ и переводящий $\alpha$ в $\beta$.

Введём гомоморфизм $\psi: F(\alpha) \to F(\beta)$, который сохраняет $F$ и 
$\psi(\alpha) = \beta$. Понятно, что в силу этих условий и сохранения
операций сложения и умножения гомоморфизм задаётся однозначно.

Понятно, что $\psi$ действует так: 

$$
\frac{f(\alpha)}{g(\alpha)} \mapsto \frac{f(\beta)}{g(\beta)}
$$

Поэтому очевидно, что $\psi$ биекция.
Так как $\beta \in F(\alpha)$, то $F(\beta) = F(\alpha)$, поэтому можно
считать, что $\psi$ -- автоморфизм.

б) Количество автоморфизмов $F(\alpha)$, сохраняющих $F$, равно количеству сопряжённых (над $F$) к $\alpha$ элементов, лежащих внутри $F(\alpha)$.

Покажем, что при автоморфизме $\alpha$ может переходить только в сопряженные
элементы.

Обозначим его минимальный многочлен $\alpha$ над $F$ за $f(x)$. 
Получаем, что $\varphi(f(\alpha)) = \varphi(0) = 0$. При этом,
так как $\varphi$ сохраняет $F$, $f(\varphi(\alpha)) = 0$.
Так как $\varphi(\alpha)$ является корнем минимального многочлена элемента
$\alpha$, то их минимальные многочлены совпадают. По определению это и
означает, что $\alpha$ и $\varphi(\alpha)$ сопряжены.

По пункту а) каждый автоморфизм, который переводит $\alpha$ в сопряженный,
который лежит внутри $F(\alpha)$, задаётся единственным образом.
